% TOP_PROBLEM: {"id": 30006981, "title": "Breuil-Mezard conjecture: equality of the Galois and automorphic sides", "category_id": 1, "category": {"id": 1, "name": "number_theory", "display_name": "Number Theory", "description": "Properties of integers, prime numbers, Diophantine equations.", "slug": "number-theory", "order_index": 1, "created_at": "2026-07-31T15:26:25.670Z"}, "status": "open"}
% FIRST_POSTED: 2026-09-27
% !TeX program = lualatex
% Compile twice: lualatex 30006981.tex
% All references and prose are embedded; no BibTeX database or external inputs.
\documentclass[11pt,a4paper]{article}
\usepackage[margin=24mm,headheight=14pt]{geometry}
\usepackage{fontspec}
\setmainfont{Latin Modern Roman}
\setsansfont{Latin Modern Sans}
\setmonofont{Latin Modern Mono}
\usepackage{amsmath,amssymb,mathtools}
\usepackage{unicode-math}
\setmathfont{Latin Modern Math}
\usepackage{microtype}
\usepackage{enumitem,xurl,needspace}
\usepackage[dvipsnames]{xcolor}
\usepackage{hyperref}
\hypersetup{unicode=true,colorlinks=true,linkcolor=MidnightBlue,urlcolor=MidnightBlue,
 pdftitle={Research Notebook: Section 30006981},pdfauthor={Research notebook},bookmarksnumbered=true}
\usepackage{bookmark}
\usepackage{tocloft}
\setlength{\cftsecnumwidth}{3em}
\cftsetpnumwidth{2.5em}
\cftsetrmarg{4em}
\usepackage{titlesec}
\titleformat{\section}{\Large\bfseries\raggedright}{\thesection}{0.7em}{}
\usepackage{fancyhdr}
\pagestyle{fancy}\fancyhf{}
\fancyhead[L]{\small\sffamily Open Problems | Research notebook}
\fancyhead[R]{\small 228 | 27 September 2026}
\fancyfoot[C]{\thepage}
\setlength{\parindent}{0pt}
\setlength{\parskip}{5pt plus 1pt}
\setlength{\emergencystretch}{3em}
\setcounter{tocdepth}{1}
\setcounter{secnumdepth}{1}
\setlist[itemize]{leftmargin=3em,itemsep=5pt,topsep=4pt,parsep=0pt}
\allowdisplaybreaks
\newcommand{\N}{\mathbb{N}}
\newcommand{\Z}{\mathbb{Z}}
\newcommand{\Q}{\mathbb{Q}}
\newcommand{\R}{\mathbb{R}}
\newcommand{\C}{\mathbb{C}}
\newcommand{\F}{\mathbb{F}}
\newcommand{\A}{\mathbb{A}}
\newcommand{\PP}{\mathbb P}
\newcommand{\E}{\mathbb{E}}
\newcommand{\eps}{\varepsilon}
\newcommand{\dd}{\,\mathrm d}
\newcommand{\sref}[2]{\hyperlink{source:#1:#2}{[S#2]}}
\newcommand{\eref}[2]{\hyperlink{extra:#1:#2}{[E#2]}}
% Turn the next switch off for a shorter reading copy. The quotations preserve
% every exactTarget field, including qualifications omitted from a short summary.
\newif\ifshowcatalogue
\showcataloguetrue
\newcommand{\cataloguescope}[1]{\ifshowcatalogue
 \paragraph{Exact scope in the supplied catalogue.}
 \begin{quote}\small #1\end{quote}\fi}

% Independently compilable extraction; original section text retained.
\numberwithin{equation}{section}
\begin{document}
\setcounter{section}{0}
% ================================================================
% problemId: 30006981
% Canonical problem ID: 30006981
% exactTarget SHA-256: 53f165c9c3a8cdde665a1d3b6cab4de066869919c66b0afda0b57b603786cd06
\clearpage
\section*{\texorpdfstring{Breuil-Mezard conjecture: equality of the Galois and automorphic sides}{Breuil-Mezard conjecture: equality of the Galois and automorphic sides}}\label{30006981}
\subsection{Definitions and mathematical statement}
\textbf{Source correction.} The catalogue writes $\sigma_{\mathrm{Gal}}=\sigma_{\mathrm{Aut}}$ as an equality of representations, but its actual source, Volpato's workshop notes, Section 1.1, Conjecture 1, writes a \emph{numerical} equality $\mu_{\mathrm{Gal}}=\mu_{\mathrm{Aut}}$. For a residual two-dimensional representation $\bar\rho$ of $G_{\Q_p}$, a Hodge type $k$, and inertial type $\tau$, the Galois quantity is a Hilbert--Samuel multiplicity of a potentially semistable deformation ring modulo a uniformizer. The automorphic expression is
\[
 \mu_{\mathrm{Aut}}(k,\tau,\bar\rho)=\sum_{n,m}a(n,m)\mu_{n,m}(\bar\rho),
\]
where $a(n,m)$ are the multiplicities of weights
$\operatorname{Sym}^n\F^2\otimes\det^m$ in the semisimplified reduction of the prescribed type lattice. Hilbert--Samuel multiplicity is the normalized leading coefficient of the local length-growth polynomial. The source's determinant and deformation conditions must be retained. The printed representation-equality target lacks definitions supporting that literal interpretation; its wording is preserved below rather than silently declared equivalent to this numerical conjecture.
\par\smallskip\textit{Formulation and concept references:} \sref{30006981}{1}.
\cataloguescope{Conjecture (Breuil-Mezard): the equality sigma\_Gal = sigma\_Aut holds, where sigma\_Gal and sigma\_Aut are the Galois and automorphic (locally algebraic) representations attached to a mod p Galois representation.}
\subsection{Short English statement}
The cited Breuil–Mézard formulation compares a Galois deformation-ring multiplicity with a weighted automorphic multiplicity, not the undefined representation equality printed in the catalogue.
% BEGIN RESEARCH ATTEMPT 228
\subsection{Research attempt}

\paragraph{Current frontier and scope audit.}
The source defect already recorded above is substantive. On p.~2 of the
workshop notes, Conjecture~1 is the numerical identity
$\mu_{\mathrm{Gal}}=\mu_{\mathrm{Aut}}$; Conjecture~2 is its framed
Hilbert--Samuel version. Neither defines the catalogue's two supposed
representations $\sigma_{\mathrm{Gal}},\sigma_{\mathrm{Aut}}$.
\sref{30006981}{1}, Section~1.1. The literal representation equality therefore
has no specified objects on which to attempt a proof. The original text
is retained, and the deductions below are explicitly about the
\emph{source-defined numerical problem}, not a silently repaired literal
representation theorem.

The numerical problem has substantial established solutions. Pa\v{s}k\=unas
proves it, with a cycle refinement, for two-dimensional representations
of $G_{\Q_p}$ with scalar endomorphisms when $p\ge5$.
\eref{30006981}{1}, Theorem~1.1. A later approach of Feng--Le Hung constructs
cycles in higher-rank generic settings and verifies the predicted
identities for sufficiently generic tame types and small Hodge--Tate
weights. Its hypotheses cannot be discarded. \eref{30006981}{2}, Abstract
and Section~1.5. The old workshop statement is not being used as a
current assertion that every numerical case remains open. The August
2026 date generated inside the HTML rendering of the 2014 Pa\v{s}k\=unas
version is not its scientific revision date.

\paragraph{Attack route.}
Patching aims to identify geometric cycles and then obtain their
multiplicities; it requires actual compatible deformation modules.
A numerical route instead tries to determine a fixed vector of weight
multiplicities from many Hodge and inertial types. We investigate the
second route: can finitely many exact tests force all others, and what
would a failed test really show? This produces a finite-obstruction
lemma, but does not manufacture the arithmetic compatibility needed to
apply it. We also test the false implication from numerical equality
back to isomorphism of representations or local rings.

\paragraph{Attempt.}
\textbf{1. Separate fixed data, type data, and the proposed unknowns.}
Fix $p$, a residual representation $\bar\rho$, and the determinant and
deformation conventions of the source. We do not change the determinant
while comparing types. Let $\mathcal W$ be the finite set of Serre
weights used there, represented by
\[
 \operatorname{Sym}^n\F^2\otimes\det^m,
       \quad 0\le n\le p-1,\quad0\le m\le p-2.
\]
For an admissible type $t=(k,\tau)$, let $a_{t,w}\in\Z_{\ge0}$ be the
semisimplified type-lattice multiplicities, and let $e_t\in\Z_{\ge0}$
be the appropriate Galois Hilbert--Samuel multiplicity. These are the
source's numerical inputs. For the algebraic diagnostic, temporarily
treat $\mu_w$ as unknown nonnegative integers and examine
\begin{equation}\label{30006981:system}
 e_t=\sum_{w\in\mathcal W}a_{t,w}\mu_w
                      \qquad\text{for every admissible }t.
\end{equation}
In the actual conjecture, the right-hand multiplicities are prescribed
by the representation-theoretic construction; finding \emph{some}
solution of \eqref{30006981:system} would not by itself identify that
prescribed vector. This distinction is essential. The weight formula
and its determinant dependence come from \sref{30006981}{1}, Section~1.1.

\textbf{2. A finite-obstruction lemma for the numerical consistency problem.}
Let $\mathcal W$ be any finite set and $a_{t,w},e_t$ any nonnegative
integers, indexed by an arbitrary family of rows $t$. Then the following
are equivalent: every finite subsystem of \eqref{30006981:system} has a
solution in $\Z_{\ge0}^{\mathcal W}$, and the entire system has such a
solution.

Here is a proof which supplies the needed boundedness instead of invoking
compactness of an unbounded integer lattice. Call a column active if it
appears positively in some row. For each active $w$, choose a row
$t(w)$ with $a_{t(w),w}\ge1$, and let $\mathcal T_0$ be this finite
collection of rows. Any nonnegative solution of those rows satisfies
\begin{equation}\label{30006981:box}
 0\le\mu_w\le
 \left\lfloor\frac{e_{t(w)}}{a_{t(w),w}}\right\rfloor.
\end{equation}
Inactive columns can be set to zero. Thus the candidates satisfying
$\mathcal T_0$ form a finite set $S$. Under finite consistency $S$ is
nonempty. If no member of $S$ satisfied all rows, choose for each member
one row that it violates. Those finitely many rows, together with
$\mathcal T_0$, would be inconsistent, a contradiction. This proves
the lemma. In the empty-active-set case finite consistency simply says
that every $e_t$ is zero, and the zero vector works.

The proof gives a precise negative-certificate principle: failure of
nonnegative-integral compatibility can always be witnessed by finitely
many types. It does not provide an a priori bound on the complexity
of the witnessing types, nor a method for computing arbitrary $e_t$.
It also does not prove that the Galois values really are compatible.
If a known representation-theoretic vector $\mu^{\mathrm{Aut}}$ is
already supplied, the actual assertion is the stronger direct test
$e_t=a_t\cdot\mu^{\mathrm{Aut}}$ for every type.

\textbf{3. Why a few consistent rows do not extrapolate.}
Integer row relations give necessary identities. If finitely many
integers $c_t$ satisfy $\sum_t c_t a_t=0$, then any solution requires
\begin{equation}\label{30006981:rowrelation}
 \sum_t c_t e_t=0.
\end{equation}
This can be checked by exact integer linear algebra, not floating-point
least squares. But neither separate positivity nor solvability one row
at a time enforces it. For example, the rows
\[
 a_1=(1,1),\quad a_2=(1,0),\quad a_3=(0,1),
 \qquad e_1=e_2=e_3=1
\]
are individually soluble over nonnegative integers, while their
combination requires simultaneously $\mu_1+\mu_2=1$ and
$\mu_1=\mu_2=1$. Even a nonnegative rational solution is not enough:
the one-row equation $2\mu=1$ has no integer solution. These examples
are tests of the proposed inference, not fabricated deformation-ring
multiplicities for an actual $\bar\rho$.

If a finite collection of rows has full column rank over $\Q$, it fixes
at most one rational vector. Checking that vector on those rows does
not imply an identity for a new row. A missing structural theorem would
have to show that every new Galois multiplicity obeys all the same
relations as the automorphic type multiplicities, and that the resulting
vector is the prescribed one. The cycle identities in
\eref{30006981}{1}, Theorem~1.1, provide exactly such substantive geometry
in their proven scope; they are not consequences of fitting a finite
linear system.

\textbf{4. Equal leading coefficients do not identify the objects.}
Let $\kappa$ be a field, and consider the one-dimensional complete local
rings
\[
 R=\kappa[[x,y]]/(xy),\qquad S=\kappa[[x,y]]/(y^2).
\]
For every integer $N\ge0$, both have
\begin{equation}\label{30006981:length}
 \operatorname{length}(R/\mathfrak m_R^{N+1})
 =\operatorname{length}(S/\mathfrak m_S^{N+1})=2N+1.
\end{equation}
For $R$ the monomial basis is $1,x,\ldots,x^N,y,\ldots,y^N$.
For $S$ it is $1,x,\ldots,x^N,y,xy,\ldots,x^{N-1}y$, with the
second list empty when $N=0$. Thus both Hilbert--Samuel multiplicities
are two, and even their entire displayed length functions agree.
But $R$ is reduced, since $(xy)=(x)\cap(y)$, whereas the nonzero element
$y\in S$ is nilpotent. The rings are not isomorphic.

There is an equally direct representation-theoretic test. Over a field
of characteristic $p$, let a generator of $C_p$ act on $\kappa^2$
either by $I_2$ or by
\[
 J=\begin{pmatrix}1&1\\0&1\end{pmatrix}.
\]
Both define representations because $(I+N)^p=I+N^p=I$ for $N^2=0$.
Their semisimplifications have two trivial constituents. Their spaces
of invariant vectors, however, have dimensions two and one respectively,
so they are not isomorphic. Consequently neither equality of
Jordan--H\"older multiplicities nor equality of local-ring multiplicities
can justify the undefined literal representation equality in the
catalogue. A specified functor and an extension-class theorem would be
additional data, not notation that can be silently inferred.

\paragraph{Stress test.}
The finite-obstruction proof used positivity to obtain
\eqref{30006981:box}; without it an unbounded integer system need not have
this property. It is not a proof that a finite selection of the existing
BM tests is universally sufficient, because the excluded candidates may
require types of unbounded complexity to detect. More importantly,
consistency does not identify the multiplicities already fixed on the
automorphic side. The elementary examples do not disprove the genuine
numerical conjecture; they disprove shortcuts that would bypass its
arithmetic content.

No global modularity hypothesis from Theorem~1 of the workshop notes
has been imported into the local conjecture without justification.
No scalar-endomorphism or $p\ge5$ hypothesis from the proven theorem
has been erased. The rings in \eqref{30006981:length} are illustrative local
rings, not claimed to be the specified potentially semistable
deformation rings. The source's warning that its notes were made during
a workshop is retained as a reason to check later precise theorems,
not as license to invent an alternative formulation.

\paragraph{Outcome.}
\textbf{Outcome: Exploratory approach with unresolved gap.}

The source-defined numerical formulation has been separated from the
undefined literal representation statement. A finite-obstruction lemma
for nonnegative-integral multiplicity systems is proved, together with
exact counterexamples to extrapolation from isolated tests and to
recovering isomorphism from numerical invariants. These are limited
structural deductions, not a new Breuil--M\'ezard theorem.

The literal target still lacks definitions of its asserted
representations. For the numerical source formulation, the missing work
is a deformation/representation-theoretic compatibility theorem in any
scope not already covered by established results. Nothing in this
attempt supplies that theorem or establishes a universal representation
isomorphism.
% END RESEARCH ATTEMPT 228
\subsection{Sources}
\begin{itemize}\raggedright
\item[\textnormal{[S1]}]\hypertarget{source:30006981:1}{} AIM workshop problem list, p-adic Representations, Modularity, and Beyond.
\url{https://aimath.org/WWN/padicmodularity/padicmodularity.pdf}.
{\small\textit{Catalogue formulation source. Consulted 24 September 2026: Section 1.1, Conjecture 1 is an equality of numerical multiplicities, not the representation equality printed in the catalogue.}}
% BEGIN ADDED REFERENCES 228
\item[\textnormal{[E1]}]\hypertarget{extra:30006981:1}{}
V. Pa\v{s}k\=unas, \emph{On the Breuil--M\'ezard conjecture},
Duke Mathematical Journal \textbf{164} (2015), no.~2, 297--359.
Consulted arXiv:1209.5205v3, 20 March 2014, Introduction and
Theorem~1.1. \url{https://arxiv.org/html/1209.5205v3}.
The theorem assumes $p\ge5$ and scalar residual endomorphisms; the
HTML-generated 2026 internal date is not a revision date.
\item[\textnormal{[E2]}]\hypertarget{extra:30006981:2}{}
T. Feng and B. Le Hung, \emph{Mirror symmetry and the Breuil--M\'ezard
Conjecture}, arXiv:2310.07006, consulted version~4, Abstract and the
revision's Section~1.5 summary. \url{https://arxiv.org/abs/2310.07006}.
The generic-parameter, tame-type and Hodge-weight qualifications are
retained. Both research sources and the workshop PDF were consulted
27 September 2026.
% END ADDED REFERENCES 228
\end{itemize}

\end{document}
